\section{前端设计：让主观质量可评分}

作者从前端设计任务入手，因为自我评估问题在这里最为突出。在没有任何干预的情况下，Claude 倾向于产出安全、可预测但视觉上平淡无奇的布局。

\begin{knowledgebox}{GAN 的启发}
生成对抗网络（GAN）的核心思想是让 generator 和 discriminator 在对抗中共同进步。作者将这一思想迁移到 LLM agent 系统中：generator 负责生成前端页面，evaluator 负责评判质量并给出具体反馈，形成迭代改进的闭环。
\end{knowledgebox}

\subsection{四维评分标准}

两个关键洞察驱动了 harness 的设计：（1）虽然审美不能完全量化，但可以通过编码设计原则来使其``可评分''；（2）将生成和评分分离可以创造驱动 generator 提升的反馈闭环。

作者定义了四个评分维度，同时写入 generator 和 evaluator 的 prompt 中：

\begin{table}[H]
\centering
\begin{tabular}{lp{9cm}l}
\toprule
\textbf{维度} & \textbf{定义} & \textbf{权重} \\
\midrule
Design Quality & 设计是否感觉像一个连贯的整体？颜色、排版、布局、图像是否创造了独特的情绪和身份感？ & 高 \\
Originality & 是否有定制化的设计决策？还是模板布局、库默认值和 AI 生成模式的堆砌？ & 高 \\
Craft & 技术执行：排版层次、间距一致性、色彩和谐、对比度。这是能力检查而非创意检查。 & 低 \\
Functionality & 独立于审美的可用性。用户能否理解界面功能、找到主要操作、完成任务？ & 低 \\
\bottomrule
\end{tabular}
\caption{四维评分标准及其权重分配}
\end{table}

\begin{importantbox}{权重设计的理由}
作者有意加大 Design Quality 和 Originality 的权重，降低 Craft 和 Functionality 的权重。原因是 Claude 在技术执行（Craft）和功能实现（Functionality）方面本身表现不错，但在设计感和原创性方面常常产出``高度通用的 AI slop''。评分标准明确惩罚这些通用模式（如``紫色渐变覆盖白色卡片''），并通过更高权重推动模型承担更多审美风险。
\end{importantbox}

\subsection{Generator-Evaluator 闭环}

整个闭环基于 \textbf{Claude Agent SDK} 构建：

\begin{enumerate}
    \item \textbf{Generator} 根据用户 prompt 生成 HTML/CSS/JS 前端页面
    \item \textbf{Evaluator} 通过 \textbf{Playwright MCP} 与活页面交互——自主导航、截屏、仔细研究实现——然后为每个维度打分并撰写详细评论
    \item 评论反馈回 generator 作为下一轮迭代的输入
    \item Generator 做出策略决策：如果分数趋势良好，\textbf{refinement}（在当前方向上优化）；如果方向不对，\textbf{pivot}（转向完全不同的审美）
    \item 重复 5--15 轮迭代
\end{enumerate}

每轮完整运行耗时较长（evaluator 需要实际操作页面），整体可达 \textbf{4 小时}。

\begin{warningbox}{Prompt 措辞的意外影响}
评分标准中的用语会以未预期的方式引导 generator 的输出。例如，包含``the best designs are museum quality''这样的短语会推动设计向特定的视觉风格收敛。这提示我们 prompt 中的隐喻和形容词本身就是强信号。
\end{warningbox}

\subsection{迭代过程中的观察}

\begin{itemize}
    \item 分数通常随迭代提升，但\textbf{不总是线性的}——作者经常发现中间轮次的产出优于最终轮次
    \item 实现复杂度倾向于逐轮增加，generator 在 evaluator 反馈下追求更ambitious 的方案
    \item 即使在第一轮（无 evaluator 反馈），有评分标准的产出已显著优于无任何 prompt 工程的 baseline
    \item Evaluator 的校准使用了 \textbf{few-shot examples}（包含详细分数拆解），以确保判断与作者偏好对齐，并减少跨迭代的分数漂移
\end{itemize}

\begin{knowledgebox}{荷兰艺术博物馆的创意跃迁}
作者用``为一家荷兰艺术博物馆创建网站''作为 prompt。前9轮产出了一个干净的深色主题着陆页——视觉上精致但在预期之内。到第10轮，generator 彻底放弃了当前方向，将网站重新构想为一个\textbf{空间体验}：用 CSS perspective 渲染的棋盘格地板 3D 房间，艺术品以自由位置悬挂在墙上，通过门廊导航而非滚动或点击切换画廊。这种创意跃迁是作者此前从未在单次生成中看到的。
\end{knowledgebox}

\subsection{本章小结}

通过将主观审美判断分解为四个可评分维度，并建立 generator-evaluator 对抗闭环，作者成功将 Claude 的前端设计质量从``技术上可用但视觉平淡''提升到了具有审美风险和创意跃迁的水平。关键在于：（1）评分标准本身就是强 prompt 信号；（2）分离生成与评估使得调优 evaluator 的``怀疑态度''远比调优 generator 的``自我批判''更可行。

